\documentclass[12pt,a4paper]{article}

\usepackage{fontspec}
\usepackage{polyglossia}
\setmainlanguage{russian}
\setotherlanguage{english}
\setmainfont{Times New Roman}

\usepackage{amsmath,amssymb,amsthm,mathtools}
\usepackage{geometry}
\usepackage{enumitem}
\usepackage{hyperref}
\usepackage{microtype}

\geometry{margin=2.5cm}
\setlength{\parindent}{1.25cm}
\setlength{\parskip}{0.35em}
\sloppy

\hypersetup{
    colorlinks=true,
    linkcolor=black,
    urlcolor=blue,
    pdftitle={Решение открытых задач Теории Единого Целого},
    pdfauthor={}
}

\newtheorem{definition}{Определение}[section]
\newtheorem{theorem}{Теорема}[section]
\newtheorem{proposition}{Предложение}[section]
\newtheorem{corollary}{Следствие}[section]
\newtheorem{remark}{Замечание}[section]

\newcommand{\Whole}{U}
\newcommand{\Parts}{\mathcal A}
\newcommand{\SubParts}{\mathcal B}
\newcommand{\Rel}{R}
\newcommand{\Values}{\mathcal V}
\newcommand{\Vstruct}{\mathbf V}
\newcommand{\State}{S}
\newcommand{\States}{\mathcal S}
\newcommand{\Disc}{\mathsf{Disc}}
\newcommand{\Time}{\mathsf{Time}}
\newcommand{\Motion}{\mathsf{Motion}}
\newcommand{\Stab}{\operatorname{Stab}}
\newcommand{\Metric}{\mathcal F}
\newcommand{\Force}{\mathfrak F}

\title{\textbf{Решение открытых задач Теории Единого Целого}\\[0.5em]
\large Рабочая формализация главы 11}
\author{}
\date{}

\begin{document}

\maketitle
\tableofcontents
\newpage

\section*{Предварительное замечание}
\addcontentsline{toc}{section}{Предварительное замечание}

В главе 11 текущей версии Теории Единого Целого были оставлены четыре открытые задачи:

\begin{enumerate}
    \item определить природу множества значений отношений \(\Values\);
    \item установить условия, при которых отношения порождают метрику;
    \item уточнить структуру реляционного времени \(T\);
    \item вывести физические величины \(m,E,p,F,L,H\) из \(\Rel,\State,T,F\).
\end{enumerate}

Настоящий документ предлагает рабочую формализацию этих задач. Он не закрывает всю физическую теорию окончательно, но задаёт минимальный математический каркас, на котором можно строить следующую версию ТЕЦ.

\section{Решение 11.1. Множество значений отношений}

В предыдущей версии множество значений отношений обозначалось через \(\Values\). Однако простого множества недостаточно. Значение отношения должно не только существовать, но и сравниваться, компоноваться и при необходимости проецироваться в числовую величину.

Поэтому вводится не просто множество, а структура:
\[
\Vstruct := (\Values,\preceq,\circ,\rho),
\]
где:
\begin{itemize}
    \item \(\Values\) --- множество значений отношений;
    \item \(\preceq\) --- порядок интенсивности, близости или вложенности отношения;
    \item \(\circ\) --- операция композиции отношений;
    \item \(\rho:\Values\to\mathbb R_{\ge0}\) --- числовая проекция отношения.
\end{itemize}

Тогда отношение имеет вид
\[
\Rel:\Parts\times\Parts\to\Vstruct.
\]

Строго говоря, значение \(\Rel(A_i,A_j)\) принадлежит носителю \(\Values\), но интерпретируется внутри всей структуры \(\Vstruct\):
\[
\Rel(A_i,A_j)\in\Values,
\qquad
\Vstruct=(\Values,\preceq,\circ,\rho).
\]

\begin{definition}[Упорядоченная композиционная реляционная структура]
Упорядоченной композиционной реляционной структурой называется четвёрка
\[
\Vstruct=(\Values,\preceq,\circ,\rho),
\]
где \(\preceq\) задаёт сравнение отношений, \(\circ\) задаёт их композицию, а \(\rho\) задаёт числовую проекцию.
\end{definition}

Числовые величины появляются только после применения проекции. Например, расстояние может быть определено как
\[
d(A_i,A_j)=\rho\bigl(\Rel(A_i,A_j)\bigr).
\]

\begin{proposition}
Отношение в ТЕЦ не обязано быть числом. Оно может быть более богатой структурой, а числовая величина возникает только как результат проекции.
\end{proposition}

\begin{proof}
Если \(\Rel(A_i,A_j)\in\Values\), а \(\Values\) является носителем структуры \(\Vstruct\), то значение отношения может содержать порядок, композиционные свойства и числовую проекцию. Число получается лишь после применения отображения \(\rho\). Следовательно, числовое расстояние является производным, а не первичным.
\end{proof}

\textbf{Итог:} \(\Vstruct\) --- это упорядоченная композиционная реляционная структура.

\section{Решение 11.2. Условия метрики}

Пусть задана метрическая проекция
\[
\Metric:\Values\to\mathbb R_{\ge0}.
\]

Тогда функция
\[
d(A_i,A_j)=\Metric\bigl(\Rel(A_i,A_j)\bigr)
\]
является метрикой, если выполнены четыре условия.

\subsection{Неотрицательность}

Для любых частей \(A_i,A_j\in\Parts\):
\[
\Metric\bigl(\Rel(A_i,A_j)\bigr)\ge0.
\]

\subsection{Тождественность}

\[
\Metric\bigl(\Rel(A_i,A_j)\bigr)=0
\Longleftrightarrow
A_i=A_j.
\]

\subsection{Симметрия метрической проекции}

\[
\Metric\bigl(\Rel(A_i,A_j)\bigr)
=
\Metric\bigl(\Rel(A_j,A_i)\bigr).
\]

Даже если само отношение направленное,
\[
\Rel(A_i,A_j)\neq\Rel(A_j,A_i),
\]
его метрическая проекция может быть симметричной:
\[
\Metric(\Rel(A_i,A_j))=\Metric(\Rel(A_j,A_i)).
\]

\subsection{Неравенство треугольника}

Для любых \(A_i,A_j,A_k\in\Parts\):
\[
\Metric\bigl(\Rel(A_i,A_k)\bigr)
\le
\Metric\bigl(\Rel(A_i,A_j)\bigr)
+
\Metric\bigl(\Rel(A_j,A_k)\bigr).
\]

\begin{theorem}[Условие метризации отношений]
Реляционная система \((\Parts,\Rel)\) порождает метрическое пространство тогда и только тогда, когда существует проекция
\[
\Metric:\Values\to\mathbb R_{\ge0},
\]
такая что функция
\[
d(A_i,A_j)=\Metric(\Rel(A_i,A_j))
\]
удовлетворяет условиям неотрицательности, тождественности, симметрии и неравенства треугольника.
\end{theorem}

\begin{proof}
Если такая проекция существует, то функция \(d\) удовлетворяет всем аксиомам метрики, следовательно, \((\Parts,d)\) является метрическим пространством. Обратно, если отношения порождают метрику, то по определению должна существовать функция, переводящая значения отношений в неотрицательные числа, удовлетворяющие аксиомам метрики. Эта функция и есть \(\Metric\).
\end{proof}

\textbf{Итог:} не всякое отношение создаёт пространство. Пространство возникает только там, где отношения допускают метрическую проекцию.

\section{Решение 11.3. Время}

В базовой версии ТЕЦ время было задано как отображение
\[
T:\States\times\States\to\mathbb R_{\ge0}.
\]

Однако этого недостаточно. Нужно различать длительность, порядок и необратимость.

\subsection{Длительность}

Вводится функция длительности:
\[
\tau:\States\times\States\to\mathbb R_{\ge0}.
\]

Она измеряет степень различия между состояниями:
\[
\tau(S_i,S_j)\ge0.
\]

В простейшем случае длительность может быть симметричной:
\[
\tau(S_i,S_j)=\tau(S_j,S_i).
\]

То есть \(\tau\) сама по себе является мерой различия, но ещё не задаёт направление времени.

\subsection{Направленное время}

Направление задаётся отношением порядка на состояниях:
\[
S_i\prec S_j.
\]

Запись
\[
S_i\prec S_j
\]
означает, что состояние \(S_j\) следует после состояния \(S_i\) в данной истории.

Поэтому реляционное время должно иметь структуру
\[
\Time := (\tau,\prec),
\]
где:
\begin{itemize}
    \item \(\tau\) показывает, сколько различия между состояниями;
    \item \(\prec\) показывает, какое состояние следует после какого.
\end{itemize}

\subsection{Необратимость}

Для описания необратимости вводится функционал сложности или энтропии:
\[
\Omega:\States\to\mathbb R.
\]

Для истории
\[
S_0\to S_1\to S_2\to\cdots
\]
может выполняться условие
\[
\Omega(S_{k+1})\ge\Omega(S_k).
\]

Если это условие является структурным законом данной истории, то обратный ход либо запрещён, либо становится маловероятным.

\begin{definition}[Реляционное время]
Реляционным временем называется тройка
\[
\Time=(\tau,\prec,\Omega),
\]
где \(\tau\) задаёт длительность, \(\prec\) задаёт порядок состояний, а \(\Omega\) задаёт критерий необратимости.
\end{definition}

\begin{proposition}
Время в ТЕЦ не сводится к числовой мере различия состояний. Оно включает различие, порядок и возможную необратимость.
\end{proposition}

\textbf{Итог:}
\[
\Time=(\tau,\prec,\Omega).
\]

Время --- это не только различие состояний, а различие, порядок и необратимость.

\section{Решение 11.4. Физические величины}

Пусть состояние системы \(\SubParts\subseteq\Parts\) на шаге \(k\) имеет вид
\[
S_k(\SubParts)=\{\Rel_k(A_i,A_j)\mid A_i,A_j\in\SubParts,\ i\neq j\}.
\]

Тогда изменение отношений между двумя соседними состояниями задаётся как
\[
\Delta\Rel(A_i,A_j)=\Rel_{k+1}(A_i,A_j)-\Rel_k(A_i,A_j),
\]
если на \(\Values\) определена операция разности или её аналог.

\subsection{Движение}

Движение определяется как изменение отношения:
\[
\Motion(A_i,A_j)=\Delta\Rel(A_i,A_j).
\]

\subsection{Расстояние}

Расстояние определяется через метрическую проекцию:
\[
d_{ij}=\Metric\bigl(\Rel(A_i,A_j)\bigr).
\]

\subsection{Временной шаг}

Временной шаг между двумя состояниями:
\[
\Delta t=T(S_k,S_{k+1}).
\]

В уточнённой формализации можно писать:
\[
\Delta t=\tau(S_k,S_{k+1}).
\]

\subsection{Скорость}

Скорость между частями \(A_i\) и \(A_j\):
\[
v_{ij}=\frac{\Delta d_{ij}}{\Delta t}.
\]

\subsection{Ускорение}

Ускорение:
\[
a_{ij}=\frac{\Delta v_{ij}}{\Delta t}.
\]

\subsection{Устойчивость}

Так как масса не является первичной, сначала вводится устойчивость части:
\[
\Stab(A_i)=\frac{1}{1+\sum_j |\Delta\Rel(A_i,A_j)|}.
\]

Чем меньше изменяются отношения части \(A_i\) с остальными частями, тем больше её устойчивость.

\subsection{Масса}

Масса определяется как функция устойчивости:
\[
m(A_i)=M\bigl(\Stab(A_i)\bigr).
\]

\subsection{Импульс}

Импульс:
\[
p_{ij}=m(A_i)v_{ij}.
\]

\subsection{Сила}

Чтобы не смешивать символ силы с метрической проекцией \(\Metric\), силу обозначим через \(\Force\):
\[
\Force_{ij}=\frac{\Delta p_{ij}}{\Delta t}.
\]

\subsection{Энергия}

Энергия определяется как мера способности менять отношения:
\[
E(\SubParts)=
\sum_{i,j}\frac12 m(A_i)v_{ij}^2
+
\sum_{i,j}\Phi\bigl(\Rel(A_i,A_j)\bigr),
\]
где \(\Phi\) --- потенциальная энергия отношений.

Кинетическая часть:
\[
K(\SubParts)=\sum_{i,j}\frac12 m(A_i)v_{ij}^2.
\]

Потенциальная часть:
\[
P(\SubParts)=\sum_{i,j}\Phi\bigl(\Rel(A_i,A_j)\bigr).
\]

Тогда
\[
E(\SubParts)=K(\SubParts)+P(\SubParts).
\]

\subsection{Лагранжиан}

Лагранжиан задаётся как разность кинетической и потенциальной частей:
\[
L=K-P.
\]

То есть
\[
L(\SubParts)=
\sum_{i,j}\frac12 m(A_i)v_{ij}^2
-
\sum_{i,j}\Phi\bigl(\Rel(A_i,A_j)\bigr).
\]

\subsection{Гамильтониан}

Гамильтониан задаётся как сумма кинетической и потенциальной частей:
\[
H=K+P.
\]

То есть
\[
H(\SubParts)=
\sum_{i,j}\frac12 m(A_i)v_{ij}^2
+
\sum_{i,j}\Phi\bigl(\Rel(A_i,A_j)\bigr).
\]

\section{Финальная формула закрытия открытых задач}

С учётом предложенных решений фундаментальная цепочка принимает вид
\[
\Whole
\to
\Parts
\to
\Disc
\to
\Rel
\to
\Vstruct
\to
\State
\to
\Time
\to
 d
\to
 v,a
\to
 m,p,\Force,E,L,H.
\]

Или словами:
\[
\boxed{\text{физика возникает не из пространства и времени,}}
\]
\[
\boxed{\text{а из устойчивых изменений отношений между частями Единого Целого.}}
\]

\section*{Итог}
\addcontentsline{toc}{section}{Итог}

Открытые задачи главы 11 получают следующую рабочую формализацию:

\begin{enumerate}
    \item \(\Values\) заменяется структурой \(\Vstruct=(\Values,\preceq,\circ,\rho)\).
    \item Пространство возникает только при существовании метрической проекции \(\Metric:\Values\to\mathbb R_{\ge0}\).
    \item Время уточняется как \(\Time=(\tau,\prec,\Omega)\).
    \item Физические величины выводятся из изменений отношений, устойчивости и метрической проекции.
\end{enumerate}

Таким образом, глава 11 перестаёт быть списком нерешённых проблем и становится программой перехода от реляционной математики к реляционной физике.

\end{document}
